% =============================================================================
% APPENDIX FM: METHOD-FRAMEWORK: Scientific Reusability vs One-Time Consulting
% =============================================================================
% Module: BCM2_00_META
% Version: 1.0 (January 2026)
% Status: Draft
%
% This appendix explains why the Framework approach differs fundamentally from
% traditional consulting and how EBF implements scientific reusability.
%
% =============================================================================

% =============================================================================
% CROSS-REFERENCE STRUCTURE
% =============================================================================
%
% CHAPTER LINKAGE (Required):
% - Primary Chapter: Chapter 1 (Introduction to EBF)
% - Secondary Chapters: Chapter 17 (Intervention Design), Chapter 24 (Applications)
%
% This appendix provides foundation for:
%
% DEPENDENCIES (what this appendix REQUIRES):
% - Appendix G: Glossary (symbol definitions)
% - Appendix UNMAPPED_BBB: CORE-WHERE (parameter sources)
%
% DEPENDENTS (what REQUIRES this appendix):
% - Appendix EEE: METHOD-DESIGN (uses framework principles)
% - Appendix GGG: METHOD-CONFIG (implements reusability)
% - Appendix FFF: METHOD-REGISTRY (stores reusable assets)
%
% =============================================================================

\begin{tcolorbox}[colback=purple!5!white, colframe=purple!75!black,
    title=Chapter Linkage]

\textbf{Primary Chapter:} Chapter 1: Introduction to EBF
\begin{itemize}[nosep]
\item Section 1.2: What is EBF? $\leftrightarrow$ This Appendix Section \ref{sec:fm-theory}
\item Section 1.4: EBF vs Traditional Methods $\leftrightarrow$ This Appendix Section \ref{sec:fm-comparison}
\end{itemize}

\textbf{Secondary Chapters:}
\begin{itemize}[nosep]
\item Chapter 17: Intervention Design --- references framework methodology
\item Chapter 24: Applications --- uses reusable model components
\end{itemize}

\textbf{Reading Guide:}
\begin{itemize}[nosep]
\item For conceptual overview: Read Chapter 1 first
\item For formal methodology: This appendix provides complete specification
\item For practical application: See Appendix EEE (METHOD-DESIGN)
\end{itemize}

\end{tcolorbox}

\chapter{METHOD-FRAMEWORK: Scientific Reusability vs One-Time Consulting}
\label{app:method-framework}


\begin{tcolorbox}[colback=blue!5!white, colframe=blue!75!black,
    title=Quick Reference: Appendix UNMAPPED_FM]

\textbf{Appendix:} FM (METHOD-FRAMEWORK)

\textbf{Core Question:} Why does Framework methodology produce superior outcomes compared to traditional consulting, and how does EBF implement this?

\textbf{Key Formula:}
\begin{equation}
V_{\text{Framework}} = \sum_{t=0}^{\infty} \delta^t \cdot N_t \cdot v_{\text{insight}} \quad \gg \quad V_{\text{Consulting}} = v_{\text{insight}}
\label{eq:fm-value}
\end{equation}

\textbf{Key Concepts:}
\begin{itemize}[nosep]
\item \textbf{Reusability Principle (UNMAPPED_FM-1)}: Assets that can be applied to $N$ problems have $N \times$ value
\item \textbf{Cumulative Knowledge (UNMAPPED_FM-2)}: Each application improves the framework
\item \textbf{Falsifiability Requirement (UNMAPPED_FM-3)}: Predictions must be testable
\end{itemize}

\textbf{Cross-References:} Appendix EEE (Design), GGG (Config), FFF (Registry), BBB (Parameters)
\end{tcolorbox}

\begin{abstract}
This appendix addresses a fundamental methodological question: Why do 80\% of consulting insights disappear after implementation, while scientific knowledge accumulates over centuries? We identify five principles that distinguish Framework methodology from traditional consulting: (1) Reusability, (2) Transparency, (3) Cumulative Growth, (4) Falsifiability, and (5) Transferability. We show how EBF implements these principles through its 10C architecture, model registry, and parameter database, creating a self-improving system where each client engagement strengthens the overall framework.
\end{abstract}


\begin{tcolorbox}[
    colback=orange!5!white,
    colframe=orange!75!black,
    title={\textbf{Appendix Scope: Ziel / In-Scope / Out-of-Scope / Constraints}},
    fonttitle=\bfseries
]

\textbf{Ziel (Objective):}
\begin{itemize}[nosep]
    \item Establish the theoretical and practical foundations for why Framework methodology produces superior, cumulative outcomes compared to traditional one-time consulting
\end{itemize}

\textbf{In-Scope:}
\begin{itemize}[nosep]
    \item The five principles of Framework methodology (UNMAPPED_FM-1 to UNMAPPED_FM-5)
    \item Comparison: Framework vs Consulting approaches
    \item The Asset Hierarchy: Insight $\to$ Model $\to$ Framework
    \item EBF's implementation of Framework principles
    \item Worked example demonstrating reusability (Assura case)
\end{itemize}

\textbf{Out-of-Scope (delegiert an andere Appendices):}
\begin{itemize}[nosep]
    \item Specific model design workflows $\rightarrow$ Appendix EEE (METHOD-DESIGN)
    \item Parameter estimation techniques $\rightarrow$ Appendix UNMAPPED_LLM (METHOD-LLMMC)
    \item Model registry structure $\rightarrow$ Appendix FFF (METHOD-REGISTRY)
    \item 10C CORE questions $\rightarrow$ Appendices AAA--IE
\end{itemize}

\textbf{Constraints (Anwendungsgrenzen):}
\begin{itemize}[nosep]
    \item Framework methodology requires upfront investment in structure
    \item Not suitable for truly unique, one-time problems (rare in practice)
    \item Requires organizational commitment to knowledge management
\end{itemize}

\textbf{Lieferobjekte (Deliverables):}
\begin{itemize}[nosep]
    \item Five axioms of Framework methodology (UNMAPPED_FM-1 to UNMAPPED_FM-5)
    \item Asset Hierarchy classification system
    \item Framework Value Formula (Equation \ref{eq:fm-value})
    \item EBF Implementation Map (Table \ref{tab:fm-ebf-implementation})
    \item Worked Example: Assura Libre Passage as reusable asset
\end{itemize}

\end{tcolorbox}


%%%%%%%%%%%%%%%%%%%%%%%%%%%%%%%%%%%%%%%%%%%%%%%%%%%%%%%%%%%%%%%%%%%%%%%%%%%%%%%
\section{The Fundamental Question}
\label{sec:fm-fundamental}
%%%%%%%%%%%%%%%%%%%%%%%%%%%%%%%%%%%%%%%%%%%%%%%%%%%%%%%%%%%%%%%%%%%%%%%%%%%%%%%

\subsection{Why This Matters}

Consider two approaches to the same business problem:

\textbf{Approach A (Traditional Consulting):}
\begin{quote}
``Based on our analysis, we recommend implementing a customer retention program with the following features...''
\end{quote}

\textbf{Approach B (Framework Methodology):}
\begin{quote}
``Your situation maps to a Status-Quo-Bias $\times$ Loss-Aversion model with parameters $\sigma_{SQ} = 0.78$ and $\lambda = 2.1$. The model predicts a 0.67\% response rate [90\% CI: 0.25\%--1.38\%]. Here is the model so you can adapt it for future decisions.''
\end{quote}

The difference is not stylistic---it is epistemological. Approach A produces an \textit{answer}. Approach B produces an \textit{asset}.

\subsection{The Core Problem}

\begin{tcolorbox}[colback=red!5!white, colframe=red!75!black,
    title=The Fundamental Question]
\textbf{Why do 80\% of consulting insights disappear after implementation, while scientific knowledge accumulates over centuries?}
\end{tcolorbox}

The answer lies in \textbf{five structural differences} between ad-hoc consulting and systematic framework methodology. This appendix formalizes these differences and shows how EBF implements them.


%%%%%%%%%%%%%%%%%%%%%%%%%%%%%%%%%%%%%%%%%%%%%%%%%%%%%%%%%%%%%%%%%%%%%%%%%%%%%%%
\section{Core Theory: The Five Principles}
\label{sec:fm-theory}
%%%%%%%%%%%%%%%%%%%%%%%%%%%%%%%%%%%%%%%%%%%%%%%%%%%%%%%%%%%%%%%%%%%%%%%%%%%%%%%

\subsection{The Insight Paradox}

Traditional consulting produces valuable insights. Yet these insights rarely compound:

\begin{itemize}
    \item \textbf{McKinsey (1980s):} ``Customer retention is 5$\times$ cheaper than acquisition''
    \item \textbf{McKinsey (2020s):} ``Customer retention is 5$\times$ cheaper than acquisition''
\end{itemize}

After 40 years and thousands of engagements, the insight remains unchanged. Compare this to physics:

\begin{itemize}
    \item \textbf{Newton (1687):} $F = ma$
    \item \textbf{Einstein (1905):} $E = mc^2$ (extends Newton to high velocities)
    \item \textbf{Quantum Mechanics (1920s):} Extends to atomic scales
\end{itemize}

Physics accumulates. Consulting does not. Why?

\subsection{The Asset Hierarchy}

We propose a three-level hierarchy of knowledge assets:

\begin{table}[htbp]
\centering
\caption{The Asset Hierarchy: From Insight to Framework}
\label{tab:fm-asset-hierarchy}
\renewcommand{\arraystretch}{1.5}
\begin{tabular}{@{}lccc@{}}
\toprule
\textbf{Level} & \textbf{Asset Type} & \textbf{Applicability} & \textbf{Value Multiple} \\
\midrule
1 & Insight & 1 problem & $1 \times$ \\
2 & Model & $N$ similar problems & $N \times$ \\
3 & Framework & $N \times M$ problem classes & $N \times M \times$ \\
\bottomrule
\end{tabular}
\end{table}

\begin{description}
    \item[Level 1: Insight] A specific answer to a specific question. ``Assura should expect 0.67\% upgrade rate.'' Valid for one problem, one time.

    \item[Level 2: Model] A reusable structure with replaceable parameters. ``Insurance upgrade decisions follow $P(\text{Upgrade}) = P(\text{Aware}) \times P(\text{Value} > \text{Cost}) \times P(\text{Act})$ with parameters that vary by context.'' Valid for any similar decision.

    \item[Level 3: Framework] A systematic way to generate models. ``The 10C Framework generates models for any behavioral question by specifying WHO, WHAT, HOW, WHEN, WHERE, AWARE, READY, STAGE, HIERARCHY, and EIT.'' Valid for any behavioral economics question.
\end{description}

\subsection{The Five Principles of Framework Methodology}

%%%%%%%%%%%%%%%%%%%%%%%%%%%%%%%%%%%%%%%%%%%%%%%%%%%%%%%%%%%%%%%%%%%%%%%%%%%%%%%
\section{Axioms and Formal Definitions}
\label{sec:fm-axioms}
%%%%%%%%%%%%%%%%%%%%%%%%%%%%%%%%%%%%%%%%%%%%%%%%%%%%%%%%%%%%%%%%%%%%%%%%%%%%%%%

\begin{tcolorbox}[colback=gray!5!white, colframe=gray!75!black,
    title=Axiom UNMAPPED_FM-1: Reusability Principle]
\textbf{Statement:} An asset $A$ that solves problem class $\mathcal{P}$ has value:
\begin{equation}
V(A) = \sum_{p \in \mathcal{P}} v(p) \cdot \mathbb{P}(\text{encounter } p)
\end{equation}
where $v(p)$ is the value of solving problem $p$.

\textbf{Interpretation:} The value of an asset scales with the number of problems it can solve. A model that applies to 100 clients is 100$\times$ more valuable than an insight for 1 client.

\textbf{Implication:} Rational knowledge investment should prioritize reusable assets over one-time solutions.
\end{tcolorbox}

\begin{tcolorbox}[colback=gray!5!white, colframe=gray!75!black,
    title=Axiom UNMAPPED_FM-2: Transparency Principle]
\textbf{Statement:} For an asset $A$ to be reusable, it must expose its assumptions $\mathcal{A}$, parameters $\Theta$, and boundary conditions $\mathcal{B}$:
\begin{equation}
A_{\text{reusable}} = \langle \mathcal{M}, \mathcal{A}, \Theta, \mathcal{B} \rangle
\end{equation}

\textbf{Interpretation:} A ``black box'' recommendation cannot be adapted. Only transparent models can be modified for new contexts.

\textbf{Implication:} Every EBF deliverable must document its model $\mathcal{M}$, assumptions $\mathcal{A}$, parameters $\Theta$, and when it does/doesn't apply $\mathcal{B}$.
\end{tcolorbox}

\begin{tcolorbox}[colback=gray!5!white, colframe=gray!75!black,
    title=Axiom UNMAPPED_FM-3: Cumulative Growth Principle]
\textbf{Statement:} Each application $t$ of a framework $F$ should update it:
\begin{equation}
F_{t+1} = \text{Update}(F_t, \text{Data}_t, \text{Feedback}_t)
\end{equation}

\textbf{Interpretation:} Frameworks learn. Each client engagement provides data that improves parameter estimates, identifies edge cases, and refines boundary conditions.

\textbf{Implication:} EBF maintains a learning loop: predictions $\to$ observations $\to$ parameter updates $\to$ improved predictions.
\end{tcolorbox}

\begin{tcolorbox}[colback=gray!5!white, colframe=gray!75!black,
    title=Axiom UNMAPPED_FM-4: Falsifiability Requirement]
\textbf{Statement:} A framework $F$ must generate testable predictions $\hat{y}$ that can be compared to observations $y$:
\begin{equation}
\exists \, y^* : |y^* - \hat{y}| > \epsilon \implies \text{Reject } F
\end{equation}

\textbf{Interpretation:} If a framework cannot be wrong, it cannot be improved. Consulting insights that are unfalsifiable (``culture matters'') cannot accumulate.

\textbf{Implication:} Every EBF model must produce quantitative predictions with confidence intervals. Deviations trigger model revision.
\end{tcolorbox}

\begin{tcolorbox}[colback=gray!5!white, colframe=gray!75!black,
    title=Axiom UNMAPPED_FM-5: Transferability Principle]
\textbf{Statement:} A framework $F$ must be transferable to new users $u$ without the original creators:
\begin{equation}
\text{Performance}(u_{\text{new}}, F) \geq (1 - \epsilon) \cdot \text{Performance}(u_{\text{creator}}, F)
\end{equation}

\textbf{Interpretation:} If knowledge dies with the consultant, it is not a framework. Frameworks must be documented, teachable, and usable by others.

\textbf{Implication:} EBF provides skills, templates, and workflows that any trained user can apply without relying on the original developers.
\end{tcolorbox}


%%%%%%%%%%%%%%%%%%%%%%%%%%%%%%%%%%%%%%%%%%%%%%%%%%%%%%%%%%%%%%%%%%%%%%%%%%%%%%%
\section{Key Results: Framework vs Consulting Comparison}
\label{sec:fm-comparison}
%%%%%%%%%%%%%%%%%%%%%%%%%%%%%%%%%%%%%%%%%%%%%%%%%%%%%%%%%%%%%%%%%%%%%%%%%%%%%%%

\begin{table}[htbp]
\centering
\caption{Structural Comparison: Traditional Consulting vs Framework Methodology}
\label{tab:fm-comparison}
\renewcommand{\arraystretch}{1.4}
\begin{tabular}{@{}p{3.5cm}p{5cm}p{5cm}@{}}
\toprule
\textbf{Dimension} & \textbf{Traditional Consulting} & \textbf{Framework Methodology} \\
\midrule
\textbf{Primary Output} & Recommendation (answer) & Model + Parameters (asset) \\
\textbf{Reusability} & One client, one time & $N$ clients, $M$ times \\
\textbf{Transparency} & ``Trust our experts'' & Full model specification \\
\textbf{Learning} & Stays in consultant's head & Documented in registry \\
\textbf{Falsifiability} & Qualitative (``it worked'') & Quantitative (0.67\% vs 0.8\%) \\
\textbf{Transferability} & Requires same team & Any trained user \\
\textbf{Pricing Model} & Per engagement & Per model + updates \\
\textbf{Client Dependency} & High (need consultants back) & Low (own the model) \\
\bottomrule
\end{tabular}
\end{table}

\subsection{The Value Equation}

The net present value of knowledge under each approach:

\textbf{Traditional Consulting:}
\begin{equation}
V_{\text{Consulting}} = v_{\text{insight}} - c_{\text{engagement}}
\end{equation}

One insight, one cost, no compounding.

\textbf{Framework Methodology:}
\begin{equation}
V_{\text{Framework}} = \sum_{t=0}^{\infty} \delta^t \cdot N_t \cdot v_{\text{insight}} - c_{\text{framework}} - \sum_{t=0}^{\infty} \delta^t \cdot c_{\text{maintain}}
\end{equation}

where:
\begin{itemize}[nosep]
    \item $\delta$ = discount rate
    \item $N_t$ = number of applications in period $t$
    \item $c_{\text{framework}}$ = upfront framework development cost
    \item $c_{\text{maintain}}$ = ongoing maintenance cost
\end{itemize}

\begin{proposition}[Framework Dominance]
For any $N > N^*$ where:
\begin{equation}
N^* = \frac{c_{\text{framework}} - c_{\text{engagement}}}{v_{\text{insight}} - c_{\text{maintain}}}
\end{equation}
the Framework approach dominates Traditional Consulting.
\end{proposition}

For EBF, with $c_{\text{framework}} \approx 50 \times c_{\text{engagement}}$ and $c_{\text{maintain}} \approx 0.1 \times v_{\text{insight}}$, we get $N^* \approx 55$ applications. After 55 uses, EBF provides infinite value relative to consulting.


%%%%%%%%%%%%%%%%%%%%%%%%%%%%%%%%%%%%%%%%%%%%%%%%%%%%%%%%%%%%%%%%%%%%%%%%%%%%%%%
\section{Integration: How EBF Implements Framework Principles}
\label{sec:fm-integration}
%%%%%%%%%%%%%%%%%%%%%%%%%%%%%%%%%%%%%%%%%%%%%%%%%%%%%%%%%%%%%%%%%%%%%%%%%%%%%%%

\begin{table}[htbp]
\centering
\caption{EBF Implementation of Framework Principles}
\label{tab:fm-ebf-implementation}
\renewcommand{\arraystretch}{1.4}
\begin{tabular}{@{}p{3cm}p{4cm}p{6cm}@{}}
\toprule
\textbf{Principle} & \textbf{EBF Component} & \textbf{How It Works} \\
\midrule
UNMAPPED_FM-1: Reusability & 10C Framework & Same 10 questions apply to ANY behavioral problem \\
UNMAPPED_FM-2: Transparency & Model Registry & Every model documents $\mathcal{M}, \mathcal{A}, \Theta, \mathcal{B}$ \\
UNMAPPED_FM-3: Cumulative & Session Database & Each engagement updates parameter estimates \\
UNMAPPED_FM-4: Falsifiability & Monte Carlo + CI & Predictions have 90\% confidence intervals \\
UNMAPPED_FM-5: Transferability & Skills \& Templates & \texttt{/design-model} usable by any trained analyst \\
\bottomrule
\end{tabular}
\end{table}

\subsection{Connection to CORE Appendices}

\begin{tcolorbox}[colback=yellow!5!white, colframe=yellow!75!black,
    title=Integration: FM $\leftrightarrow$ CORE Appendices]
\textbf{What FM provides:}
\begin{itemize}[nosep]
\item Methodological foundation for why 10C structure works
\item Justification for parameter database (BBB)
\item Framework for model registry organization (FFF)
\end{itemize}

\textbf{What FM requires:}
\begin{itemize}[nosep]
\item 10C CORE questions as the reusable structure (AAA--IE)
\item Parameter estimates from literature (BBB)
\item Model templates for reusability (EEE)
\end{itemize}
\end{tcolorbox}

\subsection{Connection to METHOD Appendices}

\begin{tcolorbox}[colback=yellow!5!white, colframe=yellow!75!black,
    title=Integration: FM $\leftrightarrow$ METHOD Appendices]
\textbf{What FM provides:}
\begin{itemize}[nosep]
\item Theoretical justification for the 9-step design workflow (EEE)
\item Rationale for default configurations (GGG)
\item Principles for model registry organization (FFF)
\end{itemize}

\textbf{What FM requires:}
\begin{itemize}[nosep]
\item Concrete implementation of design workflow (EEE)
\item Pre-configured defaults for rapid deployment (GGG)
\item Storage system for reusable models (FFF)
\end{itemize}
\end{tcolorbox}


%%%%%%%%%%%%%%%%%%%%%%%%%%%%%%%%%%%%%%%%%%%%%%%%%%%%%%%%%%%%%%%%%%%%%%%%%%%%%%%
\section{Worked Example: Assura Libre Passage}
\label{sec:fm-example}
%%%%%%%%%%%%%%%%%%%%%%%%%%%%%%%%%%%%%%%%%%%%%%%%%%%%%%%%%%%%%%%%%%%%%%%%%%%%%%%

\subsection{The Assura Case}

\paragraph{Setup.} Assura, a Swiss health insurer, offers ``Optima Plus Varia'' customers a ``Libre Passage'' option: upgrade to ``Ultra-Varia'' without health screening when hospital lists change. Price difference: +170 CHF/month. Communication: physical letter in December. Deadline: 1 month.

\paragraph{Question.} How many customers will upgrade?

\subsection{Traditional Consulting Approach}

A traditional consulting engagement would:
\begin{enumerate}
    \item Interview stakeholders (2 weeks)
    \item Conduct customer surveys (4 weeks)
    \item Analyze competitors (2 weeks)
    \item Develop recommendation (2 weeks)
\end{enumerate}

\textbf{Output:} ``Based on our analysis, we expect 1--3\% of customers to upgrade. We recommend improving the communication materials.''

\textbf{Reusability:} Zero. The next health insurer with a similar question starts from scratch.

\subsection{EBF Framework Approach}

Using the 10C Framework:

\textbf{Step 1: Model Selection}
\begin{itemize}[nosep]
    \item WHAT: Opt-in decision with financial trade-off $\to$ Prospect Theory
    \item WHEN: December mailing, 1-month deadline $\to$ Timing modifiers
    \item AWARE: Physical letter, ``not elaborate'' $\to$ Attention funnel
    \item HOW: Status quo vs. change $\to$ Status Quo Bias
\end{itemize}

\textbf{Step 2: Model Specification}
\begin{equation}
P(\text{Upgrade}) = P(\text{Aware}) \times \sum_i P(\text{Segment}_i) \times P(\text{Value}_i > \text{Cost}) \times P(\text{Act}_i) \times \tau_{\text{timing}}
\end{equation}

\textbf{Step 3: Parameter Estimation}
From BBB (CORE-WHERE) parameter database:
\begin{itemize}[nosep]
    \item $\lambda_{\text{CH}} = 2.1$ (Swiss loss aversion)
    \item $\sigma_{\text{SQ}} = 0.78$ (Status quo bias)
    \item $P(\text{Open}) = 0.92$ (Swiss postal statistics)
\end{itemize}

\textbf{Step 4: Prediction}
\begin{equation}
P(\text{Upgrade}) = 0.21 \times 0.057 \times 0.58 = 0.0069 = 0.69\%
\end{equation}

Monte Carlo simulation: \textbf{0.67\% [90\% CI: 0.25\%--1.38\%]}

\subsection{Reusability Demonstration}

The Assura model becomes a reusable asset:

\begin{table}[htbp]
\centering
\caption{Reusability: Assura Model Applied to Other Cases}
\label{tab:fm-reusability}
\renewcommand{\arraystretch}{1.3}
\begin{tabular}{@{}p{3.5cm}p{3cm}p{6cm}@{}}
\toprule
\textbf{New Client} & \textbf{Adaptation} & \textbf{What Changes} \\
\midrule
Helsana (similar product) & Parameter swap & Replace $\lambda_{\text{CH}}$ with Helsana customer profile \\
CSS (different canton) & Context update & Adjust $\Psi_{\text{cantonal}}$ modifiers \\
German insurer & Culture shift & Replace Swiss parameters with German equivalents \\
B2B insurance & Segment revision & Replace consumer segments with business decision-makers \\
\bottomrule
\end{tabular}
\end{table}

\textbf{Time for new analysis:}
\begin{itemize}[nosep]
    \item Traditional consulting: 8--10 weeks
    \item EBF with existing model: 2--4 hours
\end{itemize}

\subsection{Learning Loop}

After Assura implements the Libre Passage option:
\begin{enumerate}
    \item \textbf{Observe:} Actual upgrade rate = X\%
    \item \textbf{Compare:} Predicted 0.67\%, Observed X\%
    \item \textbf{Update:} If $|X - 0.67| > 0.2$, revise segment parameters
    \item \textbf{Store:} New case in Case Registry with observed outcome
    \item \textbf{Improve:} Future predictions benefit from calibration
\end{enumerate}


%%%%%%%%%%%%%%%%%%%%%%%%%%%%%%%%%%%%%%%%%%%%%%%%%%%%%%%%%%%%%%%%%%%%%%%%%%%%%%%
\section{Implications for Practice}
\label{sec:fm-practice}
%%%%%%%%%%%%%%%%%%%%%%%%%%%%%%%%%%%%%%%%%%%%%%%%%%%%%%%%%%%%%%%%%%%%%%%%%%%%%%%

\begin{tcolorbox}[colback=green!5!white, colframe=green!75!black,
    title=Practical Implications]
\begin{enumerate}
\item \textbf{Invest in Structure First:} Building the 10C framework took years. But now every new question is faster because the structure exists.

\item \textbf{Document Everything:} Every model, every parameter, every assumption. The Session Database (\texttt{model-building-session.yaml}) captures all decisions for future reference.

\item \textbf{Make Predictions Falsifiable:} ``0.67\% [0.25\%--1.38\%]'' can be checked. ``1--3\%'' cannot be falsified (any outcome fits).

\item \textbf{Close the Loop:} Track outcomes. Update parameters. The framework improves with use.

\item \textbf{Transfer Knowledge:} Skills like \texttt{/design-model} encode expertise. New analysts can use the framework without reinventing it.
\end{enumerate}
\end{tcolorbox}


%%%%%%%%%%%%%%%%%%%%%%%%%%%%%%%%%%%%%%%%%%%%%%%%%%%%%%%%%%%%%%%%%%%%%%%%%%%%%%%
\section{Summary}
\label{sec:fm-summary}
%%%%%%%%%%%%%%%%%%%%%%%%%%%%%%%%%%%%%%%%%%%%%%%%%%%%%%%%%%%%%%%%%%%%%%%%%%%%%%%

\begin{tcolorbox}[colback=green!10!white, colframe=green!75!black,
    title=\textbf{Appendix UNMAPPED_FM: Summary}]

\textbf{Core Question:} Why do consulting insights disappear while scientific knowledge accumulates?

\textbf{Main Contribution:}
\begin{itemize}[nosep]
    \item Five principles distinguishing Framework from Consulting (UNMAPPED_FM-1 to UNMAPPED_FM-5)
    \item Asset Hierarchy: Insight $\to$ Model $\to$ Framework
    \item Framework Value Formula showing $N^* \approx 55$ break-even point
    \item EBF implementation map for each principle
\end{itemize}

\textbf{Key Outputs:}
\begin{itemize}[nosep]
    \item Framework dominance proof: After 55 uses, infinite relative value
    \item Reusability demonstration: 8--10 weeks $\to$ 2--4 hours
\end{itemize}

\textbf{Integration:} This appendix provides the methodological foundation for all METHOD appendices (EEE, FFF, GGG) and justifies the 10C CORE structure.

\end{tcolorbox}


%%%%%%%%%%%%%%%%%%%%%%%%%%%%%%%%%%%%%%%%%%%%%%%%%%%%%%%%%%%%%%%%%%%%%%%%%%%%%%%
\section{Glossary of Symbols}
\label{sec:fm-glossary}
%%%%%%%%%%%%%%%%%%%%%%%%%%%%%%%%%%%%%%%%%%%%%%%%%%%%%%%%%%%%%%%%%%%%%%%%%%%%%%%

\begin{tcolorbox}[colback=blue!5!white, colframe=blue!75!black]
\textbf{Central Reference:} For complete symbol definitions, see
\textbf{Appendix G} (\textit{Glossary and Notation}).

This section lists only symbols introduced or specialized in this appendix.
\end{tcolorbox}

\vspace{0.5em}

\begin{table}[htbp]
\centering
\caption{Symbols Introduced in Appendix UNMAPPED_FM}
\label{tab:fm-symbols}
\renewcommand{\arraystretch}{1.3}
\begin{tabular}{@{}clp{6cm}c@{}}
\toprule
\textbf{Symbol} & \textbf{Name} & \textbf{Definition} & \textbf{Also in G?} \\
\midrule
$V_{\text{Framework}}$ & Framework Value & NPV of framework methodology & NEW \\
$V_{\text{Consulting}}$ & Consulting Value & Value of single consulting engagement & NEW \\
$N^*$ & Break-even Point & Applications needed for framework dominance & NEW \\
$\mathcal{A}$ & Assumptions & Set of model assumptions & \checkmark \\
$\mathcal{B}$ & Boundary Conditions & Where model applies/doesn't apply & NEW \\
$\mathcal{M}$ & Model & Formal model specification & \checkmark \\
$\Theta$ & Parameters & Model parameters & \checkmark \\
$\delta$ & Discount Factor & Time preference parameter & \checkmark \\
\bottomrule
\end{tabular}
\end{table}


%%%%%%%%%%%%%%%%%%%%%%%%%%%%%%%%%%%%%%%%%%%%%%%%%%%%%%%%%%%%%%%%%%%%%%%%%%%%%%%
\section{Critical Foundations and Objections}
\label{sec:fm-foundations}
%%%%%%%%%%%%%%%%%%%%%%%%%%%%%%%%%%%%%%%%%%%%%%%%%%%%%%%%%%%%%%%%%%%%%%%%%%%%%%%

\subsection{Objection 1: ``Every Problem is Unique''}

\textbf{Objection:} Business problems are unique. A framework cannot capture the nuances of each situation. Experienced consultants provide judgment that models cannot.

\textbf{Response:} Problems appear unique at the surface level but share deep structure. The Assura case looks unique (Swiss health insurance, Libre Passage option). But the underlying model---Status Quo Bias $\times$ Loss Aversion $\times$ Attention Funnel---applies to any opt-in decision with financial trade-offs. The framework captures the structure; parameters capture the uniqueness.

\textbf{Evidence:} Physics doesn't need a new theory for each falling object. $F = ma$ applies to apples, rockets, and planets. Similarly, Prospect Theory applies to insurance decisions, investment choices, and policy compliance.

\subsection{Objection 2: ``Frameworks are Too Rigid''}

\textbf{Objection:} Frameworks force problems into pre-defined boxes. Real-world complexity requires flexibility that frameworks cannot provide.

\textbf{Response:} This confuses rigidity with structure. The 10C Framework is highly flexible: it generates different models for different problems. The structure ensures completeness (all 10 dimensions considered); the parameters allow infinite variation.

\textbf{Evidence:} The Assura model and a climate policy model use the same 10C structure but have completely different parameters, segments, and predictions.

\subsection{Objection 3: ``Clients Want Answers, Not Models''}

\textbf{Objection:} Executives don't have time to understand models. They want clear recommendations, not frameworks.

\textbf{Response:} This is a communication choice, not a methodological constraint. EBF produces both: a clear recommendation (``Expect 0.67\% upgrade rate'') AND the underlying model. The model is for future use; the recommendation is for today's decision.

\textbf{Evidence:} Clients who initially wanted ``just the answer'' often return asking ``Can we use that model for our new product?'' The model's value becomes apparent over time.


%%%%%%%%%%%%%%%%%%%%%%%%%%%%%%%%%%%%%%%%%%%%%%%%%%%%%%%%%%%%%%%%%%%%%%%%%%%%%%%
\section{Open Issues and Research Agenda}
\label{sec:fm-open}
%%%%%%%%%%%%%%%%%%%%%%%%%%%%%%%%%%%%%%%%%%%%%%%%%%%%%%%%%%%%%%%%%%%%%%%%%%%%%%%

\begin{table}[htbp]
\centering
\caption{Research Roadmap for Framework Methodology}
\label{tab:fm-roadmap}
\renewcommand{\arraystretch}{1.3}
\begin{tabular}{@{}clcl@{}}
\toprule
\textbf{\#} & \textbf{Issue} & \textbf{Priority} & \textbf{Status} \\
\midrule
1 & Optimal framework granularity & HIGH & Active research \\
2 & Automated model selection & MEDIUM & Prototype in development \\
3 & Framework-to-framework transfer & MEDIUM & Theoretical stage \\
4 & Client adoption barriers & LOW & Documented but not addressed \\
\bottomrule
\end{tabular}
\end{table}

\subsection{Issue 1: Optimal Framework Granularity}

How abstract should a framework be? Too abstract (``people have preferences'') provides no predictive power. Too specific (``Swiss German-speaking male insureds aged 45--55 in December'') is not reusable. Finding the optimal level of abstraction is an open research question.

\subsection{Issue 2: Automated Model Selection}

Currently, analysts select which 10C dimensions to emphasize based on judgment. Can this be automated? Given a problem description, can the system recommend which model template to use? Early prototypes using similarity matching show promise.


%%%%%%%%%%%%%%%%%%%%%%%%%%%%%%%%%%%%%%%%%%%%%%%%%%%%%%%%%%%%%%%%%%%%%%%%%%%%%%%
\section{References}
\label{sec:fm-references}
%%%%%%%%%%%%%%%%%%%%%%%%%%%%%%%%%%%%%%%%%%%%%%%%%%%%%%%%%%%%%%%%%%%%%%%%%%%%%%%

\begin{tcolorbox}[colback=blue!5!white, colframe=blue!75!black]
\textbf{Master Bibliography:} For complete reference details, see
\texttt{bcm\_master.bib} in the repository.

This section lists appendix-specific references and data sources.
\end{tcolorbox}

\vspace{0.5em}

\subsection{Key References for This Appendix}

The following works are particularly relevant for the content of this appendix:

\begin{itemize}[nosep]
    \item \textbf{Framework Methodology:} \citet{kuhn1962structure}, \citet{lakatos1970falsification}
    \item \textbf{Reusability in Software:} \citet{parnas1972criteria}, \citet{gamma1994design}
    \item \textbf{Behavioral Economics Foundations:} \citet{kahneman1979prospectprospect}, \citet{thaler2008nudge}
    \item \textbf{Knowledge Management:} \citet{nonaka1995knowledge}, \citet{davenport1998working}
\end{itemize}

\nocite{bcm_master}


%%%%%%%%%%%%%%%%%%%%%%%%%%%%%%%%%%%%%%%%%%%%%%%%%%%%%%%%%%%%%%%%%%%%%%%%%%%%%%%
% CROSS-REFERENCE FOOTER (Required)
%%%%%%%%%%%%%%%%%%%%%%%%%%%%%%%%%%%%%%%%%%%%%%%%%%%%%%%%%%%%%%%%%%%%%%%%%%%%%%%

\vspace{1em}
\hrule
\vspace{0.5em}

\noindent\textit{Cross-references:}
Appendix G (Central Glossary),
Appendix EEE (METHOD-DESIGN),
Appendix FFF (METHOD-REGISTRY),
Appendix GGG (METHOD-CONFIG),
Appendix UNMAPPED_BBB (CORE-WHERE).

\noindent\textit{Master Bibliography:} \texttt{bcm\_master.bib}


% =============================================================================
% END OF APPENDIX FM
% =============================================================================
